% Guard: macros.tex is \input both by main.tex (after the packages it needs,
% e.g. amsmath for \DeclareMathOperator) and by individual section files so
% they remain compilable on their own. Loading it twice would redefine every
% \newcommand and abort, so the second and later loads are no-ops.
\ifdefined\PQRmacrosloaded\endinput\fi
\let\PQRmacrosloaded\relax

% --- 0. revision marking ---
% Text added to Section 2 during the current revision is wrapped in \added{...}
% so the advisors can see at a glance what is new. Redefine to \relax to print
% the paper without the marking.
\newcommand{\added}[1]{\textcolor{blue}{#1}}
\long\def\added#1{\begingroup\color{blue}#1\endgroup}  % \long: may span paragraphs


% --- A. operator-style score and function names ---
\newcommand{\HD}{\mathit{HD}}                        % histogram distance; use \HD_\gamma
\newcommand{\Attr}{\mathit{Attr}}                    % attribute set of a schema
\newcommand{\Attrs}{\mathit{Attrs}}                  % attribute set argument
% Abbreviated in MATH only; the prose still spells the names out in full.
% Spelling them out inside formulas makes displayed equations far too wide
% for a two-column layout (the product-rule step below overflowed by 200pt).
\newcommand{\Use}{\mathit{Use}}
\newcommand{\Cov}{\mathit{Cov}}
\newcommand{\Rep}{\mathit{Rep}}
\newcommand{\DivS}{\mathit{Div}}
\newcommand{\TVD}{\mathit{TVD}}
\newcommand{\GS}{\mathit{GS}}                        % global sensitivity
\DeclareMathOperator{\Dom}{Dom}
\DeclareMathOperator{\Lap}{Laplace}

% --- B. core objects ---
\newcommand{\Qone}{Q_1}
\newcommand{\Qtwo}{Q_2}
\newcommand{\pD}{p(D)}
\newcommand{\views}{\mathcal{P}}
\newcommand{\hist}[2]{H_{#1,#2}}                     % \hist{\Qone(\pD)}{A}
\newcommand{\dhist}{\Delta H}
\newcommand{\nhist}{\widetilde{\Delta H}}            % the released noisy object

% --- C. privacy budget and parameters ---
\newcommand{\eps}{\epsilon}
\newcommand{\epsP}{\epsilon_P}
\newcommand{\epsS}{\epsilon_S}
\newcommand{\epsH}{\epsilon_H}
\newcommand{\gammaref}{\gamma_{\mathrm{ref}}}
\newcommand{\gammaeval}{\gamma_{\mathrm{eval}}}
\newcommand{\taumax}{\tau_{\max}}

% --- D. mechanisms ---
\newcommand{\ExpMech}{\textsc{ExpMech}}
\newcommand{\OneShotTop}[1]{\textsc{OneShotTop-}#1}
\newcommand{\AdaptiveTau}{\textsc{AdaptiveTau}}
% --- E. experiment methods (used in sections/experiments.tex); a method is renamed here ---
% (the figure legends are drawn by src/pqde/paper_figures.py and must be redrawn after a rename)
\newcommand{\PrivExpl}{\textsf{PrivExpl}}
\newcommand{\OptPrivExpl}{\textsf{OptPrivExpl}}
\newcommand{\BF}{\textsf{BF}}
\newcommand{\SingleStage}{\textsf{SingleStage}}
\newcommand{\RandomExpl}{\textsf{Random}}
\newcommand{\Bone}{\ensuremath{\mathsf{B1}}}
\newcommand{\Btwo}{\ensuremath{\mathsf{B2}}}
\newcommand{\Bthree}{\ensuremath{\mathsf{B3}}}

\section{Extensions}
\label{sec:extension}

The analysis of \sys{} relies on a single property of the query pair, namely that both outputs are selections over $D$, so that a tuple changes each output by at most itself.
We examine two other sources of query pairs in light of this property.
For contrasts between groups the property holds and the outputs are moreover disjoint, whereas for consistent query answering it fails.

\subsection{Contrasts Between Groups}
\label{subsec:ext-contrast}

DPXPlain~\cite{TaoGMR22} explains why, in the noisy answer to a group-by query with grouping attribute $A_{gb}$ and aggregate $agg(A_{agg})$, the value of a group $g_1$ is larger than that of a group $g_2$.
We read the group-by query as one query per group, so the question becomes the pair of selections $Q_i=\sigma_{\varphi\wedge A_{gb}=g_i}$ for $i=1,2$, where $\varphi$ is the condition of the query, and an explanation shows where, and in which attribute, the two populations whose aggregate values differ are composed differently.
Both outputs are selections, so \sys{} applies as it is.
Since a tuple belongs to at most one of them, two of its bounds tighten.

\begin{proposition}[Disjoint outputs]
\label{prop:contrast-sensitivity}
Let $Q_1,Q_2$ be selections with mutually exclusive conditions and $\gamma_R=\alpha N$.
Then (1)~$\Delta_{\Use}(N)\le 2/N$ for every $N\ge1/\alpha$, (2)~the sensitivity of $\HD_\gamma$ remains $1/(1+\gamma)$, and (3)~the two histograms of a candidate 
% together 
have $L_1$-sensitivity $1$.
\end{proposition}

Intuitively, a tuple now contributes to one output only, so the view stage can be calibrated to $2/\bar N$ and the Laplace noise to scale $k/\epsilon_H$, which halves the noise of both steps.
This choice depends on the queries alone.
The model changes in two respects (proof and details in \ifpaper the full version~\cite{fullversion}\else \Cref{app:proof-contrast,app:contrast-shares}\fi).
First, $A_{gb}$ is removed from the views and from $\Attrs$, since it separates the outputs by definition, whereas $A_{agg}$ is kept, so that a view on one of its ranges relates an explanation to the aggregate.
Second, \sys{} scores how the two populations differ and not, as DPXPlain does, how much a predicate contributes to the difference of the aggregate values.
Its sensitivities therefore do not depend on the range of $A_{agg}$, but an explanation may be unrelated to the aggregate, and $\HD_\gamma$ compares counts, which differ from the shares that an \texttt{AVG} question concerns when the groups differ in size.

\subsection{Consistent Query Answering}
\label{subsec:ext-cqa}

Given denial constraints $\Sigma$ that $D$ violates, the \emph{consistent answer} $Q_\Sigma(D)$ to a selection query $Q$ consists of the tuples of $Q(D)$ that are involved in no violation~\cite{arenas1999,bertossi2011,chomicki2004}, and the pair $(Q(D),Q_\Sigma(D))$ shows which part of the selected population the inconsistencies remove.
Here the second output is not a selection, because a tuple that is added to $D$ removes from it every tuple with which it is in a violation.
Under the dependency $\texttt{AreaCode}\to\texttt{State}$, for instance, one tuple with an existing area code and a different state removes all the tuples with that area code, so the sensitivities of Proposition~\ref{prop:stage-sensitivity} no longer hold.
We see two ways to proceed (details in \ifpaper the full version~\cite{fullversion}\else \Cref{app:proof-cqa,app:cqa-details}\fi).

\paratitle{Option 1 (bounded conflicts).}
Assume a public bound $c\ge1$ on the number of tuples with which a tuple can be in a violation, as for a key constraint when at most $c+1$ tuples share a key value.

\begin{proposition}[Bounded conflicts]
\label{prop:cqa-sensitivity}
Let $Q_1=Q$, $Q_2=Q_\Sigma$ and $\gamma_R=\alpha N$.
Between neighboring databases that satisfy the bound $c$, the smaller of which has size $N\ge1/\alpha$, (1)~the usefulness of every view changes by at most $(4+c)/N$, (2)~the smoothed histogram distance of every candidate by at most $(c+1)/(1+\gamma)$, and (3)~its pair of histograms by at most $c+1$ in $L_1$ norm.
\end{proposition}

The algorithms adapt by using these scales, and Theorem~\ref{thm:framework-dp} then holds over the databases that satisfy the bound.
The noise grows linearly with $c$.
Enforcing a bound, for example by a projection to bounded degree~\cite{KasiviswanathanNRS13,DayLL16,MohapatraGHK25}, is left for future work.

\paratitle{Option 2 (conflicts as part of the input).}
The second option moves the difficulty outside \sys{}.
The detection of conflicts is treated as part of the black box, like a refinement mechanism, so that every tuple carries a flag that states whether it is involved in a violation, and $Q_2$ selects the unflagged tuples of $Q(D)$.
Both outputs are then selections and \sys{} applies unchanged, but its guarantee is with respect to neighboring \emph{annotated} databases and does not cover the effect of a tuple on the flags of the others, whose privacy cost belongs to the annotation step.
Section~\ref{sec:experiments} uses this option.
